\originalchapter{复数振幅与物质波}
\label{ch:waves}
\originalsection{干涉现象与复数振幅}
先看一个物理问题. 在Mach--Zehnder干涉仪中, 入射光被分成两条路径, 在第二个分束器重新汇合. 改变一条路径的光程, 两个出射口的亮暗会改变. 把入射减弱到单光子情形, 一次只留下一个探测事件, 重复实验的计数仍随两路的相对相位改变. 因此要解释的首先是探测统计, 不是给光子预先指定一条未被测量的轨迹.

如果只给两条路径各分配一个非负概率, 再把概率相加, 就无法描述一处增强、另一处抵消的变化. 量子描述先给各路径一个可相加的\textbf{概率振幅}, 汇合时相加, 最后取模方得到概率. 振幅既有大小又有相位; 相位记录类似波峰与波谷之间的相对位置. 复数是同时记录这两个量的方便工具. 下面介绍所需运算, 而不是先假定已经学过抽象态空间.

复数 $z=a+\ii b$ 中 $\ii^2=-1$, 由两个实数指定, 共轭为 $z^*=a-\ii b$, 模为 $|z|=\sqrt{a^2+b^2}$. 当 $z\ne0$ 时可写成 $z=r\e^{\ii\theta}$, 其中 $r=|z|>0$, $\theta$ 只确定到 $2\pi$ 的整数倍. Euler公式 $\e^{\ii\theta}=\cos\theta+\ii\sin\theta$ 把旋转写成乘法. 它由指数、正弦和余弦的绝对收敛幂级数, 将偶次和奇次项分开得到. 绝对收敛保证这次分组合法.

振幅相加与概率相加有不同结果. 两条不可区分的传播路径若有振幅 $z_1,z_2$, 合振幅为 $z_1+z_2$, 对应概率与 $|z_1+z_2|^2$ 成比例. 展开为 $|z_1|^2+|z_2|^2+2\operatorname{Re}(z_1^*z_2)$. 最后一项取决于相对相位, 是干涉项. 如果实验记录区分了两条路径, 不能直接用两个纯振幅相加来描述未读取记录后的统计, 必须把记录系统纳入, 再对它作概率求和.

\begin{example}
两个等大振幅分别为 $A$ 和 $A\e^{\ii\delta}$. 求总强度与 $\delta$ 的关系.
\end{example}
\begin{solution}
利用共轭乘积, $|A+A\e^{\ii\delta}|^2=|A|^2(2+2\cos\delta)=4|A|^2\cos^2(\delta/2)$. $\delta=0$ 时增强, $\delta=\pi$ 时抵消. 强度改变不意味着单条路径的振幅模改变; 改变的是交叉项.
\end{solution}

整体相位与相对相位要分开. 把整个态乘以同一个 $\e^{\ii\chi}$, 所有概率不变. 把叠加中的一个分量单独乘以相位, 一般会改变后续测量. 全局相位不表示新的物理态, 但局部相位差可以被干涉实验读取.

\originalsection{光子和物质波的尺度}
光电效应表明, 单个光子的能量为 $E=h\nu=\hbar\omega$, 其中 $\omega=2\pi\nu$, $\hbar=h/(2\pi)$. 真空光子满足 $E=pc$, 因而 $p=\hbar k$, $k=2\pi/\lambda$. 电子等有质量粒子也满足 de Broglie关系 $p=h/\lambda$; 在非相对论自由运动中, $E=p^2/(2m)$, 所以 $\omega(k)=\hbar k^2/(2m)$.

这些关系并不说明电子是经典介质的微小振动. 波函数是一种给出实验概率的复振幅. 经典波的场强、光子的电磁场与单个电子的概率振幅属于不同对象, 不能因为都含正弦就认作同一个物理量.

量子波长是否明显, 取决于与装置尺度的比较. 质量 $m=10^{-2}\,\mathrm{kg}$、速度 $v=10^{-1}\,\mathrm{m/s}$ 的小球, 波长约为 $6.6\times10^{-31}\,\mathrm m$. 这种尺度远小于实际对位置的控制精度. 冷原子速度很低, 波长能达到装置可分辨的尺度, 干涉便能够出现. 波长小不是量子理论失效, 而是某些实验量可用经典近似描述.

\originalsection{相速度与群速度}
单一平面波 $\e^{\ii(kx-\omega t)}$ 的等相位点满足 $kx-\omega t=\text{常数}$, 因而相速度为 $v_{\rm ph}=\omega/k$. 平面波在整条实线上都存在, 没有局域中心. 要描述局域粒子, 必须叠加不同波数; 群速度近似先考虑频谱集中在 $k_0$ 附近的情形.

令两个分量的波数为 $k_0\pm\Delta k/2$, 频率为 $\omega_0\pm\Delta\omega/2$. 相加得到
\[
2\e^{\ii(k_0x-\omega_0t)}\cos\!\left(\frac{\Delta k\,x-\Delta\omega\,t}{2}\right).
\]
快速振荡由前因子控制, 包络由余弦控制. 包络速度为 $\Delta\omega/\Delta k$; 窄频谱极限中成为 $v_g=\omega'(k_0)$. 对非相对论粒子, $v_g=\hbar k_0/m=p_0/m$, 而 $v_{\rm ph}=p_0/(2m)$. 两波拍频只是包络示意, 还不是可归一化的局域态. 窄谱波包的主包络近似以群速度传播; 二阶色散会使包络展宽. 自由非相对论粒子的平均位置速度严格为 $\langle p\rangle/m$, 一般多峰波包却未必有单一的峰值速度. 对相对论自由粒子, $\mathrm dE/\mathrm dp=pc^2/E=v$, 因而群速度仍是粒子速度.

\begin{remark}
由自由粒子色散关系和叠加原理构造运动方程, 是建立薛定谔方程的动机. 它不能仅凭经典力学推出全部量子测量规则; Born概率解释和量子态结构仍是物理理论的输入.
\end{remark}
